%% kapa: exam=1 round=37 year=2026 period=1 form=A subject=real_estate subject_ko=부동산학원론 seq=27 number=107
%% source: Q-Net 감정평가사 기출문제 (cst00309) artlSeq=5258924 file=raw/1008/5258924/2026년 제37회 감정평가사 1차 1교시.pdf
%% conversion: kordoc 4.24.51 -> tools/gen1.py (auto); review=proofread-sonnet; flags=none
\begin{question}{107}
다음 조건의 주택담보대출이 있다. 6회차에 변동금리를 적용하여 산정되는 원리금상환액에 가장 가까운 값은? (단, 주어진 조건에 한함)
\begin{framedtext}
○ 대출원금은 26,000만원이며 대출기간은 20년이다. \par
○ 대출원리금은 원리금균등분할방식으로 매년 말 1회 상환한다. \par
○ 1회차부터 5회차까지 고정금리를 적용하며, 대출금리는 연 2.4 \%이다. \par
○ 6회차부터는 변동금리가 적용되며, 6회차의 대출금리는 연 6.0 \%이다. \par
○ 원리금상환액은 각 금리적용 시점의 잔여 대출기간과 잔액을 기준으로 한다. \par
○ 연금의 현가계수(2.4 \%, 15년): 12.4729, 연금의 현가계수(2.4 \%, 20년): 15.7374 \par
○ 연금의 현가계수(6.0 \%, 15년): 9.7122, 연금의 현가계수(6.0 \%, 20년): 11.4699 \par
\end{framedtext}
\begin{choices}
  \item 1,652만원
  \item 1,796만원
  \item 2,122만원
  \item 2,267만원
  \item 2,676만원
\end{choices}
\end{question}
\answer{3}
